% =========================================================
% 9. search_read() ORM Method
% =========================================================
\section{\texttt{search\_read()} ORM Method}

\begin{frame}[standout]
  \texttt{search\_read()} ORM Method
\end{frame}

\begin{frame}[fragile]{\texttt{SEARCH\_READ()} METHOD}
  \begin{itemize}
    \item The \texttt{search\_read()} method combines \texttt{search()} and \texttt{read()}
          in one call.
    \item It finds records with a domain, then returns their field values as dictionaries.
    \item It is heavily used by the Odoo web client.
  \end{itemize}
  \begin{block}{Method Syntax}
\begin{lstlisting}
@api.model
def search_read(self, domain=None, fields=None,
                offset=0, limit=None, order=None):
    return super().search_read(domain=domain, fields=fields,
                               offset=offset, limit=limit, order=order)
\end{lstlisting}
  \end{block}
\end{frame}

\begin{frame}[fragile]{BREAKING DOWN THE METHOD SIGNATURE}
  \begin{itemize}
    \item \texttt{domain}: Filter conditions (same format as \texttt{search()}).
    \item \texttt{fields}: List of field names to return.
    \item \texttt{offset}: Number of records to skip.
    \item \texttt{limit}: Maximum number of records.
    \item \texttt{order}: Sort expression.
  \end{itemize}
\begin{lstlisting}
domain = [('active', '=', True)]
fields = ['name', 'age', 'email']
\end{lstlisting}
\end{frame}

\begin{frame}[fragile]{WHAT DOES \texttt{SEARCH\_READ()} RETURN?}
  \begin{itemize}
    \item It returns a \textbf{list of dictionaries}.
  \end{itemize}
\begin{lstlisting}
data = self.env['student.student'].search_read(
    [('age', '>=', 18)],
    ['name', 'age'],
    limit=2,
    order='age desc',
)
print(data)
\end{lstlisting}
  Output is (example):
\begin{lstlisting}
[
    {'id': 22, 'name': 'Amina', 'age': 25},
    {'id': 15, 'name': 'Ahmed', 'age': 22},
]
\end{lstlisting}
\end{frame}

\begin{frame}[fragile]{\texttt{SEARCH\_READ()} VS \texttt{SEARCH()} + \texttt{READ()}}
  Equivalent ideas:
\begin{lstlisting}
# Two calls
records = self.env['student.student'].search(domain, limit=10)
data = records.read(['name', 'age'])

# One call
data = self.env['student.student'].search_read(
    domain, ['name', 'age'], limit=10,
)
\end{lstlisting}
  \begin{itemize}
    \item \texttt{search\_read()} is convenient and often more efficient for RPC.
    \item Inside server business logic, \texttt{search()} + recordset access is usually clearer.
  \end{itemize}
\end{frame}

\begin{frame}[fragile]{WHEN TO USE \texttt{SEARCH\_READ()}}
  \begin{itemize}
    \item Building API responses.
    \item Returning table/list data to the UI.
    \item Exporting selected fields without keeping a recordset around.
  \end{itemize}
\begin{lstlisting}
rows = self.env['student.student'].search_read(
    [('gender', '=', 'female')],
    ['name', 'email', 'birth_date'],
    order='name asc',
)
\end{lstlisting}
\end{frame}

\begin{frame}[fragile]{METHOD CALL}
  \textbf{Method Call}
\begin{lstlisting}
rows = self.env['student.student'].search_read(
    [('active', '=', True)],
    fields=['name', 'age'],
    offset=0,
    limit=50,
    order='id desc',
)
\end{lstlisting}
  \begin{itemize}
    \item Here, \texttt{.search\_read(...)} is the method call.
    \item Remember: return type is a list of dicts.
  \end{itemize}
\end{frame}
